\chapter{Shallow-but-Wide: Why Brains Are Parallel, Not Deep}
\label{ch:shallow-wide}

\epigraph{%
  The brain is not a deep learning system.
  It is a massively parallel, remarkably shallow one.
  This is not a bug --- it is the architecture that makes
  real-time motor control possible.
}{}

\section{The Speed Constraint}

The most fundamental constraint on neural computation is \textbf{speed}. Every
biological computation has a time cost determined by the physics of synaptic
transmission:

\begin{neurobox}{Neural Timing Constraints}{timing}
  \begin{itemize}
    \item \textbf{Synaptic delay}: 1--5~ms per synapse (chemical synapses)
    \item \textbf{Axonal conduction}: 1--100~m/s (depends on myelination and diameter)
    \item \textbf{Action potential duration}: $\sim$1~ms
    \item \textbf{Refactory period}: $\sim$1--2~ms absolute, $\sim$3--5~ms relative
    \item \textbf{Total per-layer delay}: $\sim$5--10~ms including integration time
  \end{itemize}
\end{neurobox}

These delays are not technological limitations to be overcome with faster hardware.
They are \emph{fundamental physical constraints} of electrochemical signaling. No
amount of biological optimization can reduce synaptic delay below $\sim$1~ms.

\subsection{The Depth Budget}

Given the timing constraints, we can compute the maximum depth of any neural
computation that must complete within a given time window:

\begin{definition}{The Neural Depth Budget}{depth-budget}
  For a motor task requiring a response within time $T_{\text{response}}$, the
  maximum number of serial processing layers is:
  \begin{equation}
    L_{\max} = \frac{T_{\text{response}}}{\tau_{\text{layer}}},
    \label{eq:ch4:depth-budget}
  \end{equation}
  where $\tau_{\text{layer}} \approx 5$--$10$~ms is the per-layer processing time.
  For typical motor control scenarios:
  \begin{center}
  \renewcommand{\arraystretch}{1.3}
  \begin{tabular}{@{}lrr@{}}
  \toprule
  \textbf{Task} & $T_{\text{response}}$ & $L_{\max}$ \\
  \midrule
  Spinal reflex & 30~ms & 3--6 \\
  Postural correction & 100~ms & 10--20 \\
  Voluntary movement initiation & 150--250~ms & 15--50 \\
  Golf swing correction (mid-downswing) & 50~ms & 5--10 \\
  Saccadic eye movement & 200~ms & 20--40 \\
  \bottomrule
  \end{tabular}
  \end{center}
  Compare this with modern deep learning: GPT-4 has $\sim$120 transformer layers;
  ResNet-152 has 152 layers. These run on silicon at $\sim$ns per operation.
  Biology cannot afford this depth at $\sim$ms per operation.
\end{definition}

\section{The Motor Cortex: Six Layers, Not Sixty}

The cerebral cortex, including the primary motor cortex (M1), has a remarkably
consistent architecture across mammalian species \cite{Mountcastle1997}:

\begin{enumerate}
  \item \textbf{Layer I}: molecular layer (mostly dendrites and axons, few cell bodies)
  \item \textbf{Layer II/III}: external pyramidal layer (cortico-cortical connections)
  \item \textbf{Layer IV}: internal granular layer (thalamic input)
  \item \textbf{Layer V}: internal pyramidal layer (\emph{output layer} --- corticospinal tract)
  \item \textbf{Layer VI}: multiform layer (feedback to thalamus)
\end{enumerate}

The signal path from sensory input (Layer IV) through cortical processing to motor
output (Layer V) passes through \emph{at most 3--4 synapses within the cortex itself}.
Adding thalamic relay and spinal motor neuron connections, the total sensorimotor loop
is approximately 8--12 synapses deep.

\begin{insight}{The Architecture Is the Algorithm}{arch-is-algo}
  In deep learning, we can stack layers because silicon is fast. In biology, we cannot.
  Instead, the brain uses:
  \begin{enumerate}
    \item \textbf{Width}: $\sim$10 billion cortical neurons, each receiving $\sim$10,000
          synaptic inputs. The total computational capacity is in the \emph{connections},
          not the depth.
    \item \textbf{Recurrence}: horizontal connections within layers create recurrent
          dynamics that effectively increase computational depth through temporal
          iteration (but at the cost of time).
    \item \textbf{Specialization}: different brain regions process different aspects
          of the motor task in parallel, communicating through thick fiber bundles.
    \item \textbf{Hierarchy}: a shallow hierarchy of cortex $\to$ subcortical structures
          $\to$ spinal cord $\to$ muscles, with each level operating at different
          timescales and abstraction levels.
  \end{enumerate}
\end{insight}

\section{Computational Implications}

The shallow-but-wide architecture has profound implications for motor control theory:

\subsection{Complex Computations Through Simple Local Rules}

Each neuron performs a relatively simple computation: weighted sum of inputs, passed
through a nonlinear activation function. The power of the system comes from the
\emph{collective behavior} of billions of these simple units operating in parallel.

\begin{lstlisting}[language=Python, caption={Comparison: deep vs. wide network for motor control.}]
import numpy as np
import time

def deep_network_forward(
    x: np.ndarray,
    weights: list[np.ndarray],
    biases: list[np.ndarray],
    layer_delay_ms: float = 5.0
) -> tuple[np.ndarray, float]:
    """Forward pass through a deep network with biological timing.

    Parameters
    ----------
    x : np.ndarray
        Input vector.
    weights : list of np.ndarray
        Weight matrices for each layer.
    biases : list of np.ndarray
        Bias vectors for each layer.
    layer_delay_ms : float
        Simulated per-layer delay in milliseconds.

    Returns
    -------
    tuple
        (output, total_latency_ms)
    """
    h = x
    total_latency = 0.0
    for W, b in zip(weights, biases):
        h = np.maximum(0, W @ h + b)  # ReLU activation
        total_latency += layer_delay_ms
    return h, total_latency


def wide_network_forward(
    x: np.ndarray,
    W_hidden: np.ndarray,
    b_hidden: np.ndarray,
    W_out: np.ndarray,
    b_out: np.ndarray,
    layer_delay_ms: float = 5.0
) -> tuple[np.ndarray, float]:
    """Forward pass through a wide (2-layer) network.

    Parameters
    ----------
    x : np.ndarray
        Input vector.
    W_hidden : np.ndarray
        Hidden layer weights (wide).
    b_hidden : np.ndarray
        Hidden layer biases.
    W_out : np.ndarray
        Output layer weights.
    b_out : np.ndarray
        Output biases.
    layer_delay_ms : float
        Per-layer delay.

    Returns
    -------
    tuple
        (output, total_latency_ms)
    """
    h = np.maximum(0, W_hidden @ x + b_hidden)
    y = W_out @ h + b_out
    total_latency = 2 * layer_delay_ms  # Only 2 layers
    return y, total_latency


# Compare architectures for a motor control task
# Input: 20-dim state, Output: 10-dim muscle activations
np.random.seed(42)
n_input, n_output = 20, 10
n_params_budget = 50000  # Same total parameter budget

# Deep network: 10 layers of width 50
deep_widths = [n_input] + [50] * 9 + [n_output]
deep_W = [np.random.randn(deep_widths[i+1], deep_widths[i]) * 0.1
           for i in range(len(deep_widths)-1)]
deep_b = [np.zeros(deep_widths[i+1]) for i in range(len(deep_widths)-1)]

# Wide network: 2 layers, width 2000
wide_hidden = 2000
W_h = np.random.randn(wide_hidden, n_input) * 0.1
b_h = np.zeros(wide_hidden)
W_o = np.random.randn(n_output, wide_hidden) * 0.1
b_o = np.zeros(n_output)

x = np.random.randn(n_input)

y_deep, latency_deep = deep_network_forward(x, deep_W, deep_b)
y_wide, latency_wide = wide_network_forward(x, W_h, b_h, W_o, b_o)

print(f"Deep network (10 layers x 50 wide):")
print(f"  Latency: {latency_deep:.0f} ms")
print(f"  Parameters: {sum(W.size + b.size for W, b in zip(deep_W, deep_b)):,}")

print(f"\nWide network (2 layers x 2000 wide):")
print(f"  Latency: {latency_wide:.0f} ms")
print(f"  Parameters: {W_h.size + b_h.size + W_o.size + b_o.size:,}")

print(f"\nLatency ratio: {latency_deep/latency_wide:.1f}x faster with wide network")
\end{lstlisting}

\subsection{Hierarchical Decomposition}

The brain does not solve the full-dimensional control problem in one computation.
Instead, it decomposes the problem hierarchically:

\begin{enumerate}
  \item \textbf{High level} (prefrontal cortex, $\sim$200~ms): task goal, strategy
  \item \textbf{Mid level} (premotor/SMA, $\sim$100~ms): movement plan, sequencing
  \item \textbf{Low level} (M1 + spinal cord, $\sim$30~ms): muscle activation patterns
\end{enumerate}

Each level operates on a compressed representation of the level below:

\begin{proposition}
  In a hierarchical motor control architecture with $L$ levels, each reducing
  dimensionality by a factor $r$, the effective state dimension at level $l$ is:
  \begin{equation}
    n_l = n_0 / r^l,
    \label{eq:ch4:hierarchical-dim}
  \end{equation}
  where $n_0$ is the full state dimension. The computational cost per level
  scales as $\mathcal{O}(n_l^2)$ for a single-layer network, giving total
  cost $\sum_{l=0}^{L} n_l^2 = n_0^2 \sum_{l=0}^{L} r^{-2l}$, which converges
  for $r > 1$. The curse of dimensionality is broken by hierarchy.
\end{proposition}

\subsection{Synergies as Dimensionality Reduction}

Muscle synergies (d'Avella and Bizzi 2005 \cite{dAvellaBizzi2005}) provide the
most direct evidence for dimensionality reduction in motor control. The activation
pattern of $p$ muscles is approximated as a linear combination of $k \ll p$ synergies:
\begin{equation}
  \bm{u}(t) = \sum_{i=1}^{k} c_i(t) \bm{w}_i = W \bm{c}(t),
  \label{eq:ch4:synergy-decomposition}
\end{equation}
where $W \in \Reals^{p \times k}$ is the synergy matrix (fixed spatial patterns)
and $\bm{c}(t) \in \Reals^k$ are the time-varying synergy activations.

Experimental evidence consistently shows that $k = 4$--$6$ synergies can explain
$> 90\%$ of the variance in muscle activation patterns across a wide range of
tasks \cite{Tresch2006, Ting2007}.

\begin{lstlisting}[language=Python, caption={Muscle synergy extraction using Non-negative Matrix Factorization.}]
import numpy as np

def extract_synergies(
    emg_data: np.ndarray,
    n_synergies: int,
    max_iter: int = 500,
    tol: float = 1e-6
) -> tuple[np.ndarray, np.ndarray, float]:
    """Extract muscle synergies using NMF.

    Parameters
    ----------
    emg_data : np.ndarray, shape (n_timepoints, n_muscles)
        Processed EMG data (non-negative).
    n_synergies : int
        Number of synergies to extract.
    max_iter : int
        Maximum NMF iterations.
    tol : float
        Convergence tolerance.

    Returns
    -------
    tuple
        (synergy_weights, synergy_activations, vaf)
        W: (n_muscles, n_synergies)
        C: (n_timepoints, n_synergies)
        vaf: variance accounted for (0-1)
    """
    n_t, n_m = emg_data.shape

    # Initialize with random non-negative values
    W = np.abs(np.random.randn(n_m, n_synergies)) * 0.01
    C = np.abs(np.random.randn(n_t, n_synergies)) * 0.01

    for iteration in range(max_iter):
        # Multiplicative update rules (Lee & Seung 2001)
        # Update C
        numerator = emg_data @ W
        denominator = C @ W.T @ W + 1e-10
        C *= numerator / denominator

        # Update W
        numerator = emg_data.T @ C
        denominator = W @ C.T @ C + 1e-10
        W *= numerator / denominator

        # Check convergence
        reconstruction = C @ W.T
        residual = np.sum((emg_data - reconstruction)**2)
        total = np.sum(emg_data**2)
        vaf = 1.0 - residual / total

        if iteration > 0 and abs(vaf - prev_vaf) < tol:
            break
        prev_vaf = vaf

    return W, C, vaf


# Example: simulate EMG from 8 muscles during walking
np.random.seed(42)
n_timepoints = 200
n_muscles = 8

# Ground truth: 3 synergies generating 8-muscle activation
true_W = np.array([
    [0.8, 0.6, 0.1, 0.0, 0.2, 0.0, 0.1, 0.0],  # Syn 1: extensors
    [0.0, 0.1, 0.7, 0.9, 0.0, 0.3, 0.0, 0.1],  # Syn 2: flexors
    [0.2, 0.0, 0.1, 0.0, 0.8, 0.7, 0.6, 0.9],  # Syn 3: stabilizers
]).T  # (8 muscles, 3 synergies)

t = np.linspace(0, 2*np.pi, n_timepoints)
true_C = np.column_stack([
    np.maximum(0, np.sin(t)),
    np.maximum(0, np.sin(t + 2*np.pi/3)),
    np.maximum(0, np.sin(t + 4*np.pi/3)),
])

emg_simulated = true_C @ true_W.T + np.abs(np.random.randn(n_timepoints, n_muscles)) * 0.05

# Extract synergies
W_est, C_est, vaf = extract_synergies(emg_simulated, n_synergies=3)
print(f"Synergy extraction: {vaf:.1%} variance accounted for with 3 synergies")

# Test with different numbers
for k in range(1, 7):
    _, _, vaf_k = extract_synergies(emg_simulated, n_synergies=k)
    print(f"  k={k}: VAF = {vaf_k:.1%}")
\end{lstlisting}

\section{The Shallow-Wide Hypothesis for Motor Control}

We now formalize the central thesis of this volume:

\begin{insight}{The Shallow-Wide Hypothesis}{shallow-wide-hypothesis}
  Biological motor control systems achieve high-dimensional control through
  architectures that are:
  \begin{enumerate}
    \item \textbf{Shallow}: at most 10--20 serial processing layers (limited by
          synaptic delay)
    \item \textbf{Wide}: millions to billions of units per layer (enabled by neural
          density and parallel wiring)
    \item \textbf{Modular}: functionally specialized subnetworks process different
          aspects of the task in parallel
    \item \textbf{Hierarchical}: a small number of levels, each compressing dimensionality
    \item \textbf{Recurrent}: intra-layer connections provide temporal depth without
          adding spatial depth
  \end{enumerate}
  This architecture can approximate any smooth mapping (universal approximation theorem)
  with a single hidden layer of sufficient width. The practical advantage over deep
  networks is \textbf{latency}: a 2-layer network with 10,000 neurons per layer
  completes in 10~ms. A 100-layer network with 100 neurons per layer completes in
  500~ms---too slow for motor control.

  \textbf{For real-time motor control, width is cheap and depth is expensive.}
\end{insight}

\section*{Exercises}

\begin{enumerate}
  \item \textbf{Depth Budget.}
        For a goalkeeper diving to save a penalty kick (reaction time $\sim$200~ms),
        compute the maximum neural depth. How deep can the ``controller'' be?

  \item \textbf{Width vs. Depth.}
        Using the provided code, compare the outputs of a 10-layer $\times$ 50-wide
        network with a 2-layer $\times$ 2000-wide network on random inputs. Compare
        latencies.

  \item \textbf{Synergy Analysis.}
        Using the NMF code, determine how many synergies are needed to explain
        $> 95\%$ of variance in the simulated EMG data. Does this match the
        literature value of 4--6?

  \item \textbf{(Programming.)} Implement a motor control task (reaching) using both
        a deep and a wide network architecture. Compare tracking accuracy at different
        latency budgets.

  \item \textbf{(Research.)} Find experimental evidence for the number of cortical
        layers involved in processing a motor command from M1 to spinal motor neurons.
        How many synapses does the signal cross?
\end{enumerate}
